\section{The sign-only learner on typical samples: a nearly exponential bound}
\label{sec:typical}
Theorem~\ref{thm:signonly} holds on every training set, and its scale
$d^{\Theta(\sqrt d)}$ comes from requiring $\err_r(f)$ to stay bounded
below. For independent uniform training inputs we go much further. In
this section $f=\chi_{[d]}$ is full parity, $G\sim N(0,K)$, and, as in
Section~\ref{sec:signonly},
\[
 q_S(x)=\Pp\{f(x)G(x)>0\mid f(X_i)G(X_i)>0\ \forall i\},\qquad
 m=\E[G\mid f(X_i)G(X_i)>0\ \forall i],
\]
and $F_\theta(S)=\Pp_x\{q_S(x)\le\theta\}$, with $\theta=13/20$ and
$c_\theta$ as in Lemma~\ref{lem:meanconf}. The
main result of this section, Theorem~\ref{thm:S1}, says that for $\tanh$ at every fixed
depth and gains, $n\le\exp(10^{-3}d\log\log d/\log d)$ examples leave
$F_\theta(S)\ge3/8$ with probability at least $13/16$ over the sample. The
proof follows the skeleton of the paper proof of
\cite[Theorem~10.2]{subcritical}, with three changes. The clipped-parity
lemma, which is linear in the amplitude of the posterior mean, is replaced
by a restriction lemma that is logarithmic in it
(Section~\ref{sec:restriction}). The sharp arcsine level estimates are
replaced by a two-step regularity of the $\tanh$ dual coefficients
(Section~\ref{sec:twostep}). The far-pair analysis of the sampled tail
Gram matrix is replaced by an intrinsic-dimension matrix Chernoff bound
(Section~\ref{sec:intdim}).

\emph{Status.} Theorem~\ref{thm:S1} is proved. For $n\le n_0(d)$ (defined
in Section~\ref{sec:S1}) it follows from Theorem~\ref{thm:signonly}. For
larger $n$ the proof also uses a squared coercivity theorem with an
explicit sample budget, Theorem~\ref{thm:leancoer}, which this paper does
not prove. It is a theorem of the collection's Lean project, closed with
the standard axioms, and the collection's main subcritical theorem uses it
in its Lean proof. The rest of the proof is conventional and not
formalized (Remark~\ref{rem:S1status}).

\subsection{Parity after a random restriction}\label{sec:restriction}
For integers $j,k\ge0$ let
\[
 \agr_j(k)=\max_{\deg P\le j}\Pp_{y\in\{-1,1\}^k}\{\chi_{[k]}(y)P(y)>0\}
 =1-\err_j(\chi_{[k]}),
\]
with the error computed on the $k$-cube. Thus $\agr_j(k)=1$ if $j\ge k$,
in particular for $k=0$, and Lemma~\ref{lem:abfr}, applied on the $k$-cube
to $\chi_{[k]}$, gives
\begin{equation}\label{eq:agr}
 \agr_j(k)\ \le\ 1-2^{-k}\sum_{i\le(k-j-1)/2}\binom ki
 \ \le\ \frac12+\frac{j+1}{\sqrt{2\pi k}}\qquad(0\le j<k).
\end{equation}

\begin{lemma}[Restriction lemma]\label{lem:restriction}
For every $h:\X\to\R$, every $p\in(0,1]$, every integer $j\ge0$ and every
$b>0$,
\begin{equation}\label{eq:restriction}
 \Pp_x\{f(x)h(x)\ge b\}\ \le\ \E\,\agr_j(|\mathcal I|)
 +\frac1{b^2}\sum_{U\subseteq[d]}\widehat h(U)^2\,\Pp\{|U\cap\mathcal I|>j\},
\end{equation}
where the random set $\mathcal I\subseteq[d]$ contains each coordinate
independently with probability $p$. No margin is required off the event
$\{fh\ge b\}$, and the constant term $\widehat h(\emptyset)$ carries weight
zero.
\end{lemma}
\begin{proof}
Fix a realization $I$ of $\mathcal I$ and write $x=(y,z)$ with
$y\in\{-1,1\}^I$ and $z\in\{-1,1\}^{[d]\setminus I}$; the uniform measure
on $\X$ is the product of the uniform measures on $y$ and $z$. For fixed
$z$, the function $h_z(y)=h(y,z)$ has Walsh coefficients
$\widehat{h_z}(T)=\sum_{U:\,U\cap I=T}\widehat h(U)\chi_{U\setminus I}(z)$,
$T\subseteq I$. Split $h_z=h_z^{\rm lo}+h_z^{\rm hi}$, where
$h_z^{\rm lo}$ keeps the coefficients with $|T|\le j$. Parity restricts to
$f(y,z)=\varepsilon_z\chi_I(y)$ with
$\varepsilon_z=\chi_{[d]\setminus I}(z)\in\{\pm1\}$. Where $fh\ge b$ and
$|h_z^{\rm hi}|<b$, we have $fh_z^{\rm lo}=fh-fh_z^{\rm hi}>0$. Hence
\[
 \mathbf 1\{fh\ge b\}\le\mathbf 1\{\chi_I\,\varepsilon_zh_z^{\rm lo}>0\}
 +\mathbf 1\{|h_z^{\rm hi}|\ge b\}.
\]
Since $\varepsilon_zh_z^{\rm lo}$ has degree at most $j$ on the
$|I|$-cube, the first indicator has mean at most $\agr_j(|I|)$ over $y$.
By Markov's inequality the second has mean at most
$b^{-2}\sum_{T\subseteq I,\,|T|>j}\widehat{h_z}(T)^2$ over $y$. Two
distinct sets $U,U'$ with $U\cap I=U'\cap I$ have
$U\setminus I\ne U'\setminus I$, so the characters $\chi_{U\setminus I}$
in the sum defining $\widehat{h_z}(T)$ are orthonormal in $z$, and
$\E_z\widehat{h_z}(T)^2=\sum_{U\cap I=T}\widehat h(U)^2$. Therefore
$\Pp_x\{fh\ge b\}\le\agr_j(|I|)+b^{-2}\sum_U\widehat h(U)^2\,
\mathbf 1\{|U\cap I|>j\}$. Average over $\mathcal I$.
\end{proof}

\begin{corollary}[Explicit form]\label{cor:restriction}
Let $1\le u\le d$, let $h$ have degree at most $u$, let $b>0$, and put
$\rho_u=u/d$ and
\[
 \mathcal R=\frac1{b^2}\sum_{U\ne\emptyset}\widehat h(U)^2 .
\]
For every integer $J$ with $1\le J\le e^2u$,
\begin{equation}\label{eq:restriction-explicit}
 \Pp_x\{fh\ge b\}\ \le\ \frac12+e\sqrt{\frac{J\rho_u}\pi}
 +\frac12\exp\Bigl(-\frac J{8e^2\rho_u}\Bigr)+\mathcal R\,e^{-J}.
\end{equation}
In particular, if $\varepsilon\in(0,1]$ and
$1+\log_+(\mathcal R/\varepsilon)\le e^2u$, then
\[
 \Pp_x\{fh\ge b\}\le\frac12+\varepsilon
 +e\sqrt{\frac{\rho_u\,(1+\log_+(\mathcal R/\varepsilon))}\pi}
 +\frac12e^{-1/(8e^2\rho_u)} .
\]
\end{corollary}
\begin{proof}
Take $j=J-1$ and $p=J/(e^2u)\in(0,1]$ in \eqref{eq:restriction}. Then
$|\mathcal I|\sim\mathrm{Bin}(d,p)$ has mean $dp=J/(e^2\rho_u)$. By
\eqref{eq:agr}, $\agr_j(k)\le\frac12+J/\sqrt{2\pi k}$ for $k\ge1$ (trivially
if $j\ge k$), and $\agr_j\le1$ always. Split according to whether
$|\mathcal I|\ge dp/2$ and use the Chernoff bound
$\Pp\{|\mathcal I|<dp/2\}\le e^{-dp/8}$:
\[
 \E\,\agr_j(|\mathcal I|)\le\frac12+\frac J{\sqrt{\pi dp}}
 +\frac12e^{-dp/8}
 =\frac12+e\sqrt{\frac{J\rho_u}\pi}+\frac12e^{-J/(8e^2\rho_u)}.
\]
In the second term of \eqref{eq:restriction}, $\widehat h(U)=0$ for
$|U|>u$, the weight vanishes for $U=\emptyset$, and for $1\le|U|\le u$ a
union bound over $J$-subsets of $U$ gives
$\Pp\{|U\cap\mathcal I|\ge J\}\le\binom uJp^J\le(eup/J)^J=e^{-J}$ (the
probability is zero if $J>u$). For the last claim take
$J=\max\{1,\lceil\log(\mathcal R/\varepsilon)\rceil\}
\le1+\log_+(\mathcal R/\varepsilon)$, so that
$\mathcal Re^{-J}\le\varepsilon$, and use $J\ge1$ in the exponential.
\end{proof}

With $p=1$ and $\deg h\le j$, Lemma~\ref{lem:restriction} reduces to
\eqref{eq:agr} on the whole cube. Remark~\ref{rem:loglinear} compares
Corollary~\ref{cor:restriction} with the clipped-parity lemma.

\subsection{Two-step regularity of the dual coefficients}\label{sec:twostep}
\begin{lemma}[First layer]\label{lem:twostep}
Let $\sigma>0$, $\varpi=\pi/(2\sigma)$ and
$\upsilon_m=\frac12(\sqrt{\varpi^2+4m}-\varpi)$. The dual coefficients of
$\tanh$ at variance $\sigma^2$ satisfy, for every odd $m\ge1$,
\[
 \eta(\sigma):=\frac1{6(1+\varpi)^4}\ \le\ \frac{\upsilon_m^4}{(m+1)(m+2)}
 \ \le\ \frac{c_{m+2}(\sigma^2)}{c_m(\sigma^2)}\ \le\ \frac{m^2}{(m+1)(m+2)} .
\]
\end{lemma}
\begin{proof}
By Lemma~\ref{lem:tanh},
$c_{m+2}/c_m=\sigma^4(I_{m+2}/I_m)^2/((m+1)(m+2))
=(\sigma^2\E_\mu\omega^2)^2/((m+1)(m+2))$, where $\mu$ is the probability
measure on $(0,\infty)$ with density proportional to
$g(\omega)=\omega^me^{-\sigma^2\omega^2/2}/\sinh(\pi\omega/2)$. Since
$\omega g(\omega)\sim\frac2\pi\omega^m\to0$ at $0$ and $g$ has Gaussian
decay, integration by parts gives $\int_0^\infty\omega g'(\omega)\,d\omega
=-\int_0^\infty g(\omega)\,d\omega$. As
$\omega g'/g=m-\sigma^2\omega^2-\frac{\pi\omega}2\coth\frac{\pi\omega}2$,
this reads
\[
 \sigma^2\E_\mu\omega^2=m+1-\E_\mu\Bigl[\frac{\pi\omega}2
 \coth\frac{\pi\omega}2\Bigr].
\]
Since $y\coth y\ge1$, $\sigma^2\E_\mu\omega^2\le m$, which gives the upper
bound. Since $y\coth y=y+2y/(e^{2y}-1)\le1+y$, Jensen's inequality gives
$\sigma^2\E_\mu\omega^2\ge m-\frac\pi2(\E_\mu\omega^2)^{1/2}$. With
$Y=\sigma(\E_\mu\omega^2)^{1/2}$ this is $Y^2+\varpi Y-m\ge0$, so
$Y\ge\upsilon_m=2m/(\varpi+\sqrt{\varpi^2+4m})\ge m/(\varpi+\sqrt m)$, and
$c_{m+2}/c_m=Y^4/((m+1)(m+2))\ge\upsilon_m^4/((m+1)(m+2))$. Finally, for
$m\ge1$, $\varpi+\sqrt m\le(1+\varpi)\sqrt m$ and $(m+1)(m+2)\le6m^2$, so
$\upsilon_m^4/((m+1)(m+2))\ge m^4/((1+\varpi)^4m^2\cdot6m^2)=\eta(\sigma)$.
\end{proof}

\begin{lemma}[Composites]\label{lem:composite-reg}
Let $H(t)=\sum_{m\ge1}h_mt^m$ have nonnegative coefficients with
$h_{m+2}\ge\eta h_m$ for all $m\ge1$, and let $g_0,g_1,\ldots\ge0$. Then
the coefficients of $\sum_kg_kH^k$ satisfy $[t^{m+2}]\ge\eta\,[t^m]$ for
all $m\ge1$. Consequently, for $\phi=\tanh$, every depth and all gains,
with $\eta=\eta(\sigma_1)$ and $c_1=\prod_{\ell=1}^{L-1}c_1(q_\ell)>0$ the
linear coefficient of $F$,
\begin{equation}\label{eq:twostep}
 c_{m+2}\ge\eta\,c_m\qquad\text{and}\qquad c_m\ge\eta^{(m-1)/2}c_1
 \qquad(m\text{ odd}).
\end{equation}
\end{lemma}
\begin{proof}
For $k\ge1$, $[t^m]H^k=\sum h_{j_1}\cdots h_{j_k}$ over compositions
$j_1+\cdots+j_k=m$ with all $j_i\ge1$. The map
$(j_1,\ldots,j_k)\mapsto(j_1+2,j_2,\ldots,j_k)$ injects compositions of $m$
into compositions of $m+2$, and $h_{j_1+2}\ge\eta h_{j_1}$; hence
$[t^{m+2}]H^k\ge\eta[t^m]H^k$. The term $k=0$ contributes only in degree
$0$. Multiply by $g_k$ and sum. For $\tanh$ take $H=\widehat\phi_{q_1}$,
which satisfies the hypothesis by Lemma~\ref{lem:twostep} for odd $m$ and
trivially for even $m$, and let $g_k$ be the coefficients of
$\widehat\phi_{q_{L-1}}\circ\cdots\circ\widehat\phi_{q_2}$ (the identity if
$L=2$). The linear coefficient of a composition of series without constant
term is the product of the linear coefficients, and $c_1(q)>0$ (proof of
Corollary~\ref{cor:composite}). Iterating from $m=1$ gives the second bound.
\end{proof}

The constant $\eta(\sigma)$ is crude. Numerically, the smallest ratio
$c_{m+2}(\sigma^2)/c_m(\sigma^2)$ over odd $m\le4001$ is attained at $m=1$
and equals $0.0034$, $0.0152$, $0.0601$, $0.1181$, $0.1410$ for
$\sigma=0.3,0.5,1,2,3$, against $\eta(\sigma)\approx1.1\cdot10^{-4}$,
$5.7\cdot10^{-4}$, $3.8\cdot10^{-3}$, $1.6\cdot10^{-2}$, $3.1\cdot10^{-2}$
(verification record, \texttt{verification/revision-2026-10-20/} next to
this paper).

\paragraph{Pair masses.} For $0\le j\le D$ put
$W_j=\binom Dja_j=\sum_mc_m\Pp\{\zeta_m=j\}$, where $\zeta_m$ is the
distance from its start of the walk of Lemma~\ref{lem:spectrum} after $m$
steps (proof of Lemma~\ref{lem:levels}(c)). Then
$\sum_jW_j=\sum_mc_m=1$. If $F$ is odd, $W_j=0$ for even $j$, and for odd
$r$ the pairing of Lemma~\ref{lem:spectrum}(e) and
$\binom d{j-1}+\binom dj=\binom Dj$ give
\begin{equation}\label{eq:pairs}
 \tau_r=A\sum_{j\ge r+2,\ j\text{ odd}}W_j,\qquad
 \lambda_{r+1}=\frac{A\,W_{r+2}}{\binom D{r+2}} .
\end{equation}
(If $d$ is even, the unpaired top level $q=d$ fits \eqref{eq:pairs},
because $\binom dd=1=\binom DD$.)

\begin{lemma}[Cube transfer]\label{lem:transfer}
Let $F$ be odd with $c_{m+2}\ge\eta c_m$ for all odd $m\ge1$.
\begin{enumerate}
\item[(a)] For odd $j$ with $j+2\le D$,
$W_{j+2}\ge\eta\,\frac{(D-j)(D-j-1)}{D^2}\,W_j$.
\item[(b)] For odd $r$ with $r+3\le D/2$, $\tau_r\le(1+4/\eta)\tau_{r+2}$.
\item[(c)] For odd $r$ and an integer $k\ge0$ with $r+2k+1\le D/2$,
$\tau_r\le(1+4/\eta)^k\tau_{r+2k}$.
\end{enumerate}
\end{lemma}
\begin{proof}
(a) Two outward steps give
$\Pp\{\zeta_{m+2}=j+2\}\ge\Pp\{\zeta_m=j\}\frac{D-j}D\frac{D-j-1}D$. Hence
$W_{j+2}\ge\sum_{m\text{ odd}}c_{m+2}\Pp\{\zeta_{m+2}=j+2\}
\ge\eta\frac{(D-j)(D-j-1)}{D^2}\sum_mc_m\Pp\{\zeta_m=j\}$.
(b) By \eqref{eq:pairs}, $\tau_r=AW_{r+2}+\tau_{r+2}$ and
$\tau_{r+2}\ge AW_{r+4}$. By (a),
$W_{r+4}\ge\eta\frac{(D-r-2)(D-r-3)}{D^2}W_{r+2}\ge\frac\eta4W_{r+2}$,
since $D-r-3\ge D/2$. So $AW_{r+2}\le\frac4\eta\tau_{r+2}$.
(c) Apply (b) at $r,r+2,\ldots,r+2k-2$.
\end{proof}

Numerically the ratio $\tau_r/\tau_{r+2}$ stays far below $1+4/\eta$, but
it depends strongly on the gains. Over odd $r$ with $r+3\le D/2$ and
$D\in\{41,101,201,301\}$, its largest value is between $1.6$ and $1.8$ for
$L=2$, $\sigma_1=8$ (where $1+4/\eta\approx50$), and between $270$ and
$320$ for $L=2$, $\sigma_1=0.2$ (where $1+4/\eta\approx1.5\cdot10^5$); the
other configurations in the verification record lie in between.

\subsection{The sampled tail Gram matrix}\label{sec:intdim}
For a nonzero positive semidefinite matrix $\Sigma$ write
$\operatorname{intdim}(\Sigma)=\operatorname{tr}\Sigma/\lambda_{\max}(\Sigma)\ge1$.

\begin{proposition}[An intrinsic-dimension expectation bound]\label{prop:intdim}
Let $Z_1,\ldots,Z_n$ be independent random real symmetric matrices of a
common size with $0\preceq Z_k\preceq\ell I$ almost surely, where $\ell>0$.
Let $Y=\sum_kZ_k$, and let $\Sigma\succeq\E Y$ be positive semidefinite and
nonzero. Then for every $\vartheta>0$,
\[
 \E\lambda_{\max}(Y)\ \le\ \frac{e^\vartheta-1}\vartheta\,\lambda_{\max}(\Sigma)
 +\frac\ell\vartheta\log\operatorname{intdim}(\Sigma).
\]
\end{proposition}
\begin{proof}
We use Lieb's concavity theorem \cite[Theorem~6]{lieb}: for every fixed
real symmetric $H$, the map $B\mapsto\operatorname{tr}\exp(H+\log B)$ is
concave on positive definite matrices. Conditionally on
$Z_1,\ldots,Z_{n-1}$, apply it with $H=\frac\vartheta\ell\sum_{k<n}Z_k$ and
Jensen's inequality to $B=e^{\vartheta Z_n/\ell}$; this replaces
$\frac\vartheta\ell Z_n$ by $\log\E e^{\vartheta Z_n/\ell}$.
Iterating over $k$ gives the subadditivity of matrix cumulant generating
functions \cite[Lemma~3.4]{tropp2012},
\[
 \E\operatorname{tr}\exp\Bigl(\frac\vartheta\ell Y\Bigr)\le
 \operatorname{tr}\exp\Bigl(\sum_k\log\E e^{\vartheta Z_k/\ell}\Bigr).
\]
Put $g=(e^\vartheta-1)/\ell$. Convexity gives $e^{\vartheta x/\ell}\le1+gx$
on $[0,\ell]$, so $e^{\vartheta Z_k/\ell}\preceq I+gZ_k$ and
$\E e^{\vartheta Z_k/\ell}\preceq I+g\E Z_k$. The
logarithm is operator monotone and $\log(1+x)\le x$, so
$\log\E e^{\vartheta Z_k/\ell}\preceq g\E Z_k$. As $\operatorname{tr}\exp$ is monotone
under $\preceq$, $\E\operatorname{tr}e^{\vartheta Y/\ell}\le\operatorname{tr}e^{g\E Y}
\le\operatorname{tr}e^{g\Sigma}$. Since $Y\succeq0$,
$e^{\vartheta\lambda_{\max}(Y)/\ell}-1\le\operatorname{tr}(e^{\vartheta Y/\ell}-I)$,
and Jensen's inequality gives
$e^{\vartheta\E\lambda_{\max}(Y)/\ell}-1\le\operatorname{tr}(e^{g\Sigma}-I)$.
Let $\mu_{\max}=\lambda_{\max}(\Sigma)$. The function $x\mapsto e^{gx}-1$ is
convex and vanishes at $0$, so $e^{g\mu}-1\le(\mu/\mu_{\max})(e^{g\mu_{\max}}-1)$
for every eigenvalue $\mu\in[0,\mu_{\max}]$ of $\Sigma$; summing,
$\operatorname{tr}(e^{g\Sigma}-I)\le\operatorname{intdim}(\Sigma)
(e^{g\mu_{\max}}-1)$. Since $\operatorname{intdim}(\Sigma)\ge1$, this gives
$e^{\vartheta\E\lambda_{\max}(Y)/\ell}\le\operatorname{intdim}(\Sigma)e^{g\mu_{\max}}$,
that is,
$\frac\vartheta\ell\E\lambda_{\max}(Y)\le g\mu_{\max}+\log\operatorname{intdim}(\Sigma)$.
Multiply by $\ell/\vartheta$.
\end{proof}

This is a variant of the expectation bound in the intrinsic-dimension
matrix Chernoff inequality \cite[Chapter~7]{tropp2015}; we prove it to
have explicit constants.

\begin{lemma}[Tail Gram matrix]\label{lem:tailgram}
Let $K$ be as in Section~\ref{sec:signonly}, let $0\le r<d$, and let
$X_1,\ldots,X_n$ be independent and uniform on $\X$ (repeats allowed). The
sampled matrix $Q_{r,SS}=(Q_r(X_i,X_j))_{i,j\le n}$ of the tail kernel
$Q_r=K_{>r}$ satisfies
\[
 \E\|Q_{r,SS}\|_{\rm op}\ \le\ (e-1)\,n\lambda_{r+1}
 +\tau_r\log\frac{\tau_r}{\lambda_{r+1}} .
\]
\end{lemma}
\begin{proof}
Let $\psi(x)=(\lambda_{|U|}^{1/2}\chi_U(x))_{|U|>r}$, so that
$Q_r(x,y)=\langle\psi(x),\psi(y)\rangle$. Then $Q_{r,SS}=\Psi\Psi^{\mathsf T}$,
where $\Psi$ has rows $\psi(X_i)^{\mathsf T}$, and
$\|Q_{r,SS}\|_{\rm op}=\lambda_{\max}(\Psi^{\mathsf T}\Psi)=\lambda_{\max}(Y)$
with $Y=\sum_iZ_i$ and $Z_i=\psi(X_i)\psi(X_i)^{\mathsf T}$. The $Z_i$ are
independent and positive semidefinite, with
$\lambda_{\max}(Z_i)=\|\psi(X_i)\|^2=\tau_r$. Since
$\E\chi_U\chi_V=\mathbf 1_{U=V}$, $\E Y=\Sigma:=n\operatorname{diag}
(\lambda_{|U|})_{|U|>r}$, with $\lambda_{\max}(\Sigma)=n\lambda_{r+1}$ by
monotonicity and $\operatorname{intdim}(\Sigma)=\tau_r/\lambda_{r+1}$.
Apply Proposition~\ref{prop:intdim} with $\ell=\tau_r$ and $\vartheta=1$.
\end{proof}

The lemma holds for every $n$ and involves no far-pair, near-pair or
collision analysis; compare Lemma~\ref{lem:gram}. By \eqref{eq:gram-lower}
below, $\|Q_{r,SS}\|_{\rm op}\ge\max\{n\lambda_{r+1},\tau_r\}$ for every
sample, so the bound is sharp in expectation up to the logarithmic term.

\begin{lemma}[Intrinsic dimension of the $\tanh$ tail]\label{lem:intdim-tanh}
Let $F$ be odd and satisfy \eqref{eq:twostep}, and let $r$ be odd with
$r+2\le D$. Then
\[
 \binom D{r+2}\le\frac{\tau_r}{\lambda_{r+1}}\le
 \frac{D^{r+2}}{(r+2)!\,c_{r+2}},\qquad\text{hence}\qquad
 \log\frac{\tau_r}{\lambda_{r+1}}\le(r+2)\log\frac{eD}{r+2}
 +\frac{r+1}2\log\frac1\eta+\log\frac1{c_1}.
\]
\end{lemma}
\begin{proof}
By \eqref{eq:pairs},
$\tau_r/\lambda_{r+1}=\binom D{r+2}\sum_{j\ge r+2}W_j/W_{r+2}$ (odd $j$),
and $W_{r+2}\le\sum_{j\ge r+2}W_j\le\sum_jW_j=1$. The walk moves away from
its start in each of its first $r+2$ steps with probability
$\prod_{i=0}^{r+1}\frac{D-i}D=(r+2)!\binom D{r+2}/D^{r+2}$, so
$W_{r+2}\ge c_{r+2}(r+2)!\binom D{r+2}/D^{r+2}$. Finally
$(r+2)!\ge((r+2)/e)^{r+2}$, and $c_{r+2}\ge\eta^{(r+1)/2}c_1$ by
\eqref{eq:twostep}.
\end{proof}

Corollary~\ref{cor:composite} gives instead
$\log(1/c_{r+2})\le C_\phi\sqrt{r+2}+\log(1/\kappa)$. Either way the
leading term is $\log\binom D{r+2}$, which the lower bound shows to be
sharp; at linear $r$ it is of order $d$.

\subsection{The posterior mean at the training inputs}\label{sec:mean-train}
For a training sequence $S$ and an integer $r$, let $\mathrm{Coer}(r)$ be
the event
\[
 \frac1n\sum_{i=1}^n\bigl(f(X_i)P(X_i)\bigr)_-^2\ \ge\ \frac14\|P\|_2^2
 \qquad\text{for every }P:\X\to\R\text{ with }\deg P\le r,
\]
where $t_-=\max\{-t,0\}$. For integers $d\ge1$ and $r\ge0$ and real
$x\ge1$ put
\begin{equation}\label{eq:leanbudget}
 \mathsf B(d,r,x)=2^{76}D_r\bigl(\log(8x)+\lceil\log_2D_r\rceil+1\bigr).
\end{equation}

\begin{theorem}[Squared coercivity with an explicit budget]\label{thm:leancoer}
Let $d,n\ge1$ and $r\ge0$ be integers with $64r\le d$ and
$\mathsf B(d,r,n)\le n$. If $X_1,\ldots,X_n$ are independent and uniform
on $\X$ and $f=\chi_{[d]}$, then $\Pp\{\mathrm{Coer}(r)\}\ge15/16$.
\end{theorem}

This paper does not prove Theorem~\ref{thm:leancoer}. It is the theorem
\texttt{parity\_\allowbreak squared\_\allowbreak coercivity\_\allowbreak
probability} in the file
\texttt{lean/\allowbreak NeuralGeneralization/\allowbreak
ParitySquaredSampling.lean} of the collection's Lean project
\cite{leanproject}, with the
same hypotheses $0<d$, $64r\le d$, $0<n$ and $\mathsf B(d,r,n)\le n$. In
that statement $D_r$ is \texttt{lowDimension d r}, the number of subsets of
$[d]$ of size at most $r$; $\lceil\log_2D_r\rceil$ is
\texttt{Nat.clog 2} of it; the probability is the exact product law of $n$
independent uniform points of the cube; $P$ ranges over all real Walsh
expansions whose coefficients vanish above degree $r$; and $\|P\|_2^2$ is
the uniform average of $P^2$. The constant $2^{76}$ is the one proved in
Lean and is not optimized. The Lean proof of the collection's main
subcritical theorem (\texttt{MainProofs.\allowbreak subcriticalExponential})
uses this theorem. The project's verification record
(\texttt{lean/\allowbreak VERIFICATION.md}) reports that this proof depends
only on the axioms \texttt{propext}, \texttt{Classical.choice} and
\texttt{Quot.sound}, and the log of its aggregate audit \texttt{Audit.lean}
reports the same for all 7,202 proved project declarations. The record
does not name Theorem~\ref{thm:leancoer}. But \texttt{Audit.lean} checks
every theorem of the project in the import closure of the library's root
module, and the recorded closure contains
\texttt{ParitySquaredSampling.lean} (see \texttt{Audit.log} and
\texttt{project-\allowbreak import-\allowbreak closure.json} in
\texttt{lean/\allowbreak review/\allowbreak verification/\allowbreak
robust-2026-10-18/}), so the audit covers Theorem~\ref{thm:leancoer}. Its
Lean proof is a finite chaining argument formalized in the project; it does
not use \cite[Lemma~4.1]{subcritical}. We did not rerun Lean for this
paper. The record lists the SHA-256 hashes of the 621 files of the source
snapshot audited on 18 October 2026, and we checked that each of these
files is byte-identical in the current project; the record also states
that the library sources were unchanged at the independent rebuild of
19 October. The project's only \texttt{.lean} files outside that snapshot
are a supplementary check and two exact-type audit scripts kept with the
verification evidence; the library does not import them.
For comparison, \cite[Lemma~4.4]{subcritical} gives
$\Pp\{\mathrm{Coer}(r)\}\ge1-\delta$ for each $\delta>0$ when $r\le d/64$
and $n\ge A_\delta D_r\log(2D_r)$, with a constant $A_\delta\ge1$ that is
not explicit; its proof rests on \cite[Lemma~4.1]{subcritical}
(Remark~\ref{rem:S1status}). Section~\ref{sec:obstruction} uses it for
general $\delta$.

\begin{lemma}[Inverting the budget]\label{lem:invert}
Let $d,v\ge1$ be integers, and put $a=2^{76}D_v$ and
$\ell=\lceil\log_2D_v\rceil+1$, so that $\mathsf B(d,v,x)=a(\log(8x)+\ell)$.
\begin{enumerate}
\item[(a)] The function $x\mapsto\mathsf B(d,v,x)/x$ decreases on
$[1,\infty)$. Hence, if $x_0\ge1$ and $\mathsf B(d,v,x_0)\le x_0$, then
$\mathsf B(d,v,x)\le x$ for every $x\ge x_0$.
\item[(b)] If $x\ge1$ and $x<\mathsf B(d,v,x)$, then
$x<a\,(2\log(8a)+3\ell)$.
\item[(c)] If moreover $v\le(d+1)/3$, then
$x<2^{77}\binom dv\bigl(120+7v\log\frac{ed}v\bigr)$.
\end{enumerate}
\end{lemma}
\begin{proof}
Since $D_v\ge d+1\ge2$, we have $\ell\ge2$. (a) The derivative of
$(\log(8x)+\ell)/x$ is $(1-\log(8x)-\ell)/x^2<0$ for $x\ge1$. (b) The
function $g(y)=y/(\log(8y)+\ell)$ increases on $[1,\infty)$, since
$g'(y)=(\log(8y)+\ell-1)/(\log(8y)+\ell)^2>0$. Put $N=a(2\log(8a)+3\ell)\ge1$
and $w=\log(8a)+2\ell$. Then $2\log(8a)+3\ell\le2w$, and $\log(2w)\le w$
because $e^w\ge2w$ for every real $w$. Hence
\[
 \log(8N)+\ell=\log(8a)+\log(2\log(8a)+3\ell)+\ell
 \le\log(8a)+w+\ell=2\log(8a)+3\ell,
\]
that is, $g(N)\ge a$. The hypothesis says $g(x)<a$, so $x<N$.
(c) For $1\le q\le(d+1)/3$ we have
$\binom d{q-1}=\frac q{d-q+1}\binom dq\le\frac12\binom dq$, so
$D_v\le2\binom dv$; and $\binom dv\le(ed/v)^v$. With $h=v\log(ed/v)$ this
gives $\log D_v\le\log2+h$ and $\ell\le\log_2D_v+2\le3+h/\log2$. Hence
$2\log(8a)+3\ell=158\log2+2\log D_v+3\ell\le160\log2+9+(2+3/\log2)h
<120+7h$, and $a\le2^{77}\binom dv$. Apply (b).
\end{proof}

For a sequence $S$ with $s$ distinct inputs $S'=(x'_1,\ldots,x'_s)$ and for
$0\le v<d$, write $Q_v=K_{>v}$ and $Q_{v,SS}=(Q_v(X_i,X_j))_{i,j\le n}$,
repeats included.

\begin{lemma}[The posterior mean at the training inputs]\label{lem:meantrain}
Let $K$ be as in Section~\ref{sec:signonly}. For every training sequence
$S$ of length $n\ge1$ and every $0\le v<d$:
\begin{enumerate}
\item[(a)] $m=\sum_{k\le s}w_kK(\cdot,x'_k)$ with $w=K_{S'S'}^{-1}m_{S'}$, and
$\|m\|_K^2=w^{\mathsf T}K_{S'S'}w\le s$;
\item[(b)] $\frac1n\sum_{i=1}^nm_{>v}(X_i)^2\le\frac sn\|Q_{v,SS}\|_{\rm op}$;
\item[(c)] on $\mathrm{Coer}(v)$,
$\|m_{\le v}\|_2^2\le4\frac sn\|Q_{v,SS}\|_{\rm op}\le4\|Q_{v,SS}\|_{\rm op}$.
\end{enumerate}
\end{lemma}
\begin{proof}
(a) The conditioning event $\mathcal A_S=\{f(X_i)G(X_i)>0\ \forall i\}$
depends only on $G_{S'}$, so Gaussian regression gives
$m(x)=\E[\E[G(x)\mid G_{S'}]\mid\mathcal A_S]=K_{xS'}K_{S'S'}^{-1}m_{S'}$.
Then $\|m\|_K^2=w^{\mathsf T}K_{S'\X}K^{-1}K_{\X S'}w=w^{\mathsf T}K_{S'S'}w
=m_{S'}^{\mathsf T}K_{S'S'}^{-1}m_{S'}$. The vector
$\xi=K_{S'S'}^{-1/2}G_{S'}$ is standard Gaussian, $\mathcal A_S$ is a cone
in $\xi$, and the radius of $\xi$ is independent of its direction; so
$\E[\|\xi\|^2\mid\mathcal A_S]=s$, and Jensen's inequality gives
$m_{S'}^{\mathsf T}K_{S'S'}^{-1}m_{S'}=\|\E[\xi\mid\mathcal A_S]\|^2\le s$.
(b) Projecting (a) onto degrees above $v$ gives
$m_{>v}=\sum_kw_kQ_v(\cdot,x'_k)$. Let $\nu_k$ be the multiplicity of
$x'_k$ in $S$, and define $\tilde w\in\R^n$ by $\tilde w_i=w_k/\nu_k$ when
$X_i=x'_k$. Then $(Q_{v,SS}\tilde w)_i=m_{>v}(X_i)$ and
$\tilde w^{\mathsf T}Q_{v,SS}\tilde w=w^{\mathsf T}Q_{v,S'S'}w\le
w^{\mathsf T}K_{S'S'}w\le s$, because $K=K_{\le v}+Q_v$ with both terms
positive semidefinite. Hence $\sum_im_{>v}(X_i)^2=\tilde w^{\mathsf T}
Q_{v,SS}^2\tilde w\le\|Q_{v,SS}\|_{\rm op}\,\tilde w^{\mathsf T}Q_{v,SS}
\tilde w\le s\|Q_{v,SS}\|_{\rm op}$.
(c) Put $P=m_{\le v}$. At every training input $f(X_i)m(X_i)>0$, because
$f(X_i)G(X_i)>0$ on $\mathcal A_S$ and $\Pp(\mathcal A_S)>0$. Hence
$(f(X_i)P(X_i))_-\le|m(X_i)-P(X_i)|=|m_{>v}(X_i)|$: if $fP<0<fm$, then
$|m-P|=fm-fP>-fP$. Now apply $\mathrm{Coer}(v)$ and (b).
\end{proof}

\begin{proposition}[Master inequality]\label{prop:master}
Let $K$ be as in Section~\ref{sec:signonly}. Let $0\le r\le u<d$ with
$u\ge1$ and $1\le J\le e^2u$ be integers, and let $S$ have length $n\ge1$
and lie in $\mathrm{Coer}(r)$. Then
\[
 F_\theta(S)\ \ge\ \frac12-e\sqrt{\frac{Ju}{\pi d}}
 -\frac12e^{-Jd/(8e^2u)}
 -\frac{8e^{-J}}{c_\theta^2\tau_u}\bigl(4\|Q_{r,SS}\|_{\rm op}+n\lambda_{r+1}\bigr)
 -\Bigl(2+\frac8{c_\theta^2}\Bigr)\frac{n\lambda_{u+1}}{\tau_u}.
\]
\end{proposition}
\begin{proof}
Let $h=m_{\le u}$, $b=\frac{c_\theta}2\sqrt{\tau_u/2}$,
$E_1=\{\sigma_x^2<\tau_u/2\}$ and $E_2=\{|m-h|\ge b\}$, with $\sigma_x^2$ as
in Lemma~\ref{lem:meanconf}. By Lemma~\ref{lem:meanconf}(iii), (iv) and
Markov's inequality, $\Pp(E_1)\le2n\lambda_{u+1}/\tau_u$ and
$\Pp(E_2)\le n\lambda_{u+1}/b^2=8n\lambda_{u+1}/(c_\theta^2\tau_u)$. Off
$E_1\cup E_2$, $q_S(x)>\theta$ implies $f(x)m(x)>c_\theta\sigma_x\ge2b$ by
Lemma~\ref{lem:meanconf}(ii), hence $f(x)h(x)>b$. So
$\{q_S>\theta\}\subseteq\{fh\ge b\}\cup E_1\cup E_2$. The function $h$ has
degree at most $u$, and by Lemma~\ref{lem:meantrain}(c) and
Lemma~\ref{lem:meanconf}(iv) with $s\le n$,
\[
 \mathcal R=\frac1{b^2}\sum_{U\ne\emptyset}\widehat h(U)^2
 \le\frac{\|m_{\le r}\|_2^2+\|m_{\le u}-m_{\le r}\|_2^2}{b^2}
 \le\frac8{c_\theta^2\tau_u}\bigl(4\|Q_{r,SS}\|_{\rm op}+n\lambda_{r+1}\bigr).
\]
Corollary~\ref{cor:restriction} bounds $\Pp_x\{fh\ge b\}$, and
$F_\theta(S)=1-\Pp_x\{q_S>\theta\}$.
\end{proof}

\subsection{The theorem}\label{sec:S1}
Let $\phi=\tanh$, put $c_*=10^{-3}$, $t=\lfloor\sqrt d/5\rfloor$,
$k_d=\lceil\log d\rceil$,
\[
 n_0(d)=\frac{\binom dt}{23C_0},\qquad
 n_{\max}(d)=\exp\Bigl(c_*\frac{d\log\log d}{\log d}\Bigr),
\]
\[
 \varepsilon_0=\frac1{64(2+8/c_\theta^2)},\qquad
 \mathsf B_*=32\bigl((e-1)2^{77}+1\bigr)+2^{77},
\]
with $C_0$ as in Theorem~\ref{thm:signonly}. Let $\eta=\eta(\sigma_1)$
and $c_1$ be as in Lemma~\ref{lem:composite-reg}. For odd $r$ put
\[
 \Lambda_r=120+7(r+2)\log\frac{ed}{r+2},\qquad
 \Lambda'_r=(r+2)\log\frac{eD}{r+2}+\frac{r+1}2\log\frac1\eta
 +\log\frac1{c_1},
\]
\[
 \mathsf B_r=32\bigl((e-1)2^{77}\Lambda_r+\Lambda'_r\bigr)+2^{77}\Lambda_r,
\]
and let
\[
 \mathsf T_d=\log\Bigl(1+\frac4\eta\Bigr)\Bigl(1+\frac{\log d
 +\log(2^{77}/\varepsilon_0)}{2\log\frac{\log d}{1.3c_*}}\Bigr),
\]
\[
 \mathsf J_d=\log d+1+\log\frac{256\mathsf B_*}{c_\theta^2}+\mathsf T_d,\qquad
 \mathsf r_d=1.1\Bigl(\frac{c_*}{\log d}+\frac2d\Bigr).
\]
Let $d_0\ge3000$ be such that the following hold for every $d\ge d_0$,
with the budget $\mathsf B$ of \eqref{eq:leanbudget}.
\begin{enumerate}
\item[(L1)] $\mathsf B(d,t-3,n_0(d))\le n_0(d)$.
\item[(L2)] $c_*\frac d{\log d}+2\le\min\{\frac d{64},\frac{d+1}3\}$ and
$\lceil c_*d/\log d\rceil\le d/e$.
\item[(L3)] $2k_d\le0.1(t-2)$.
\item[(L4)] $1.1(c_*\frac d{\log d}+2)\le0.1d$.
\item[(L5)] $1.25(\frac{c_*}{\log d}+\frac2d)\le\frac{1.3c_*}{\log d}$ and
$\log\frac{\log d}{1.3c_*}\ge1$.
\item[(L6)] $\Lambda_r\le d$ and $\Lambda'_r\le d$ for every odd
$r<c_*d/\log d$.
\item[(L7)] $2^{77}/\varepsilon_0\le d$.
\item[(L8)] $\mathsf J_d\le e^2(t-4)$, $\mathsf J_d\mathsf r_d\le
\pi/(256e^2)$ and $\mathsf r_d\le1/(8e^2\log32)$.
\end{enumerate}
Such a $d_0$ exists. For (L1), let $d\ge3000$, so that $t\ge10$ and
$n_0(d)\ge(5\sqrt d)^{10}/(23C_0)\ge1$. By the estimate
$\binom d{q-1}\le\frac12\binom dq$ for $q\le(d+1)/3$ (proof of
Lemma~\ref{lem:invert}) and by
$\binom d{t-3}\le\binom dt(\frac t{d-t+1})^3\le\frac8{125d^{3/2}}\binom dt$,
we have $D_{t-3}\le\frac{16}{125d^{3/2}}\binom dt$, and also
$\log(8n_0(d))+\lceil\log_2D_{t-3}\rceil+1\le6+3t\log\frac{ed}t$. Hence
\[
 \frac{\mathsf B(d,t-3,n_0(d))}{n_0(d)}\ \le\
 \frac{2^{80}\cdot23C_0}{125\,d^{3/2}}\Bigl(6+3t\log\frac{ed}t\Bigr)
 =O\Bigl(\frac{\log d}d\Bigr).
\]
Further, $\Lambda_r$ and $\Lambda'_r$ are $O(c_*d\log\log d/\log d)$ for
$r<c_*d/\log d$; $\mathsf T_d=(1+o(1))\frac12\log(1+\frac4\eta)
\frac{\log d}{\log\log d}$, so $\mathsf J_d=(1+o(1))\log d$; and
$\mathsf r_d=(1+o(1))1.1c_*/\log d$. Hence
$\mathsf J_d\mathsf r_d\to1.1c_*<\pi/(256e^2)=1.66\ldots\cdot10^{-3}$, and
every other condition holds for large $d$ as well. Every constant in
(L1)--(L8) is explicit, so $d_0$ can be computed from $\sigma_1$ and $c_1$;
it is astronomically large (Remark~\ref{rem:d0}).

\begin{theorem}[Sign-only learner, typical samples]\label{thm:S1}
Let $\phi=\tanh$, and fix $L\ge2$ and gains $\sigma_1,\ldots,\sigma_L>0$.
There is $d_0$, depending only on $\sigma_1$ and
$c_1=\prod_{\ell<L}c_1(q_\ell)$, such that for every
$d\ge d_0$ and every integer $0\le n\le n_{\max}(d)
=\exp(10^{-3}d\log\log d/\log d)$ the following holds. If
$X_1,\ldots,X_n$ are independent and uniform on $\X$ and $f=\chi_{[d]}$,
then
\[
 \Pp_S\Bigl\{\Pp_x\{q_S(x)\le\tfrac{13}{20}\}\ge\tfrac38\Bigr\}\ \ge\
 \frac{13}{16} ,
\]
and on the same event the sign-only network learner satisfies
$\liminf_{N\to\infty}F_N(S)\ge3/8$. If $n\le n_0(d)$, the inner inequality
holds for every training sequence.
\end{theorem}
The theorem is proved. Its coercivity input, Theorem~\ref{thm:leancoer},
is a theorem of the collection's Lean project, closed with the standard
axioms; the rest of the proof is conventional and not formalized
(Remark~\ref{rem:S1status}). Through $c_1$, the threshold $d_0$ depends on
the depth and on $\sigma_1,\ldots,\sigma_{L-1}$, but not on $\sigma_L$;
indeed $q_S$ does not change when $K$ is multiplied by a constant. The
limit $N\to\infty$ is taken at fixed $d$, $L$, gains and $S$. By
Remark~\ref{rem:readout}, the conclusion $\liminf_NF_N(S)\ge3/8$ also
holds, on the same event and with the same status, for the margin learner
whose readout variance is $\sigma_L^2N^{-b}$ with $b<1$: its correct-label
probabilities converge to the same $q_S(x)$. Two lemmas come first.

\begin{lemma}[Adaptive cutoff]\label{lem:cutoff}
Let $d\ge d_0$ and $n_0(d)<n\le n_{\max}(d)$.
\begin{enumerate}
\item[(a)] Some odd $r\ge1$ has $\mathsf B(d,r,n)\le n$; let $r$ be the
largest one.
\item[(b)] $t-4\le r<c_*d/\log d$; in particular $64r\le d$ and
$r+2\le(d+1)/3$.
\item[(c)] $n/\binom D{r+2}<2^{77}\Lambda_r$ and
$n\lambda_{r+1}/\tau_r<2^{77}\Lambda_r$.
\item[(d)] Put $\hat\rho=1.25(r+2)/d$ and
$k^*=\lceil\log(2^{77}\Lambda_r/\varepsilon_0)/(2\log(1/\hat\rho))\rceil$.
Then $0\le k^*\le k_d$, and $u=r+2k^*$ satisfies $u+1\le D/2$,
$\rho_u=u/d\le1.1(r+2)/d\le\mathsf r_d$ and
$n\lambda_{u+1}/\tau_u\le n/\binom D{u+2}<\varepsilon_0$.
\end{enumerate}
\end{lemma}
\begin{proof}
The budget $\mathsf B(d,v,x)$ is nondecreasing in $v$, and
$\mathsf B(d,v,x)\ge D_v$ for $x\ge1$. We also use
$\binom D{j+2}/\binom Dj=\frac{(D-j)(D-j-1)}{(j+1)(j+2)}\ge
(\frac{D-j-1}{j+2})^2$.

(a) We have $n>n_0(d)\ge1$ and $t-3\ge1$, so (L1) and
Lemma~\ref{lem:invert}(a) give $\mathsf B(d,1,n)\le\mathsf B(d,t-3,n)\le n$.

(b) Suppose $r\ge r_1:=\lceil c_*d/\log d\rceil$. Then
$n\ge\mathsf B(d,r,n)\ge D_r\ge\binom d{r_1}\ge(d/r_1)^{r_1}$. The function
$x\mapsto x\log(d/x)$ increases on $(0,d/e]$, and
$c_*d/\log d\le r_1\le d/e$ by (L2), so
$\log n\ge\frac{c_*d}{\log d}\log\frac{\log d}{c_*}>\log n_{\max}(d)$, a
contradiction. Hence $r<c_*d/\log d$, and (L2) gives the two consequences.
Now suppose $r+2\le t-3$. Maximality of $r$, monotonicity in $v$, (L1) and
Lemma~\ref{lem:invert}(a) give
$n<\mathsf B(d,r+2,n)\le\mathsf B(d,t-3,n)\le n$, a contradiction. Hence
$r\ge t-4$.

(c) By maximality $n<\mathsf B(d,r+2,n)$, and $r+2\le(d+1)/3$ by (b).
Lemma~\ref{lem:invert}(c) with $v=r+2$ gives
$n<2^{77}\binom d{r+2}\Lambda_r\le2^{77}\binom D{r+2}\Lambda_r$. The second
claim follows from Lemma~\ref{lem:levels}(b). Since
$(r+2)\log\frac{ed}{r+2}\le3rH_r$ with $H_r=\log(eD/r)>5$ for
$1\le r\le d/64$, this is a bound of the form
$n/\binom D{r+2}\le C\,rH_r$, with $\Lambda_r\le45\,rH_r$.

(d) By (L5), $\hat\rho\le1.3c_*/\log d$ and $\log(1/\hat\rho)\ge1$; by
(L6) and (L7), $\log(2^{77}\Lambda_r/\varepsilon_0)\le2\log d$. So
$k^*\le k_d$, and $k^*\ge0$ because $2^{77}\Lambda_r>\varepsilon_0$. For
odd $j\in[r+2,u]$, (L3) and $r\ge t-4$ give $j+2\le r+2k_d+2\le1.1(r+2)$,
and (L4) gives $D-j-1\ge d-1.1(r+2)\ge0.9d$. Hence
$(j+2)/(D-j-1)\le\frac{1.1}{0.9}\frac{r+2}d<\hat\rho$ and
$\binom D{j+2}\ge\hat\rho^{-2}\binom Dj$. Iterating,
$\binom D{u+2}\ge\hat\rho^{-2k^*}\binom D{r+2}$, so by (c)
$n/\binom D{u+2}<2^{77}\Lambda_r\hat\rho^{2k^*}\le\varepsilon_0$ by the
choice of $k^*$. Lemma~\ref{lem:levels}(b) at $u$ gives
$n\lambda_{u+1}/\tau_u\le n/\binom D{u+2}$. Finally
$u+1<u+2\le1.1(r+2)\le D/2$, and $\rho_u\le1.1(r+2)/d\le\mathsf r_d$ by
(b).
\end{proof}

\begin{lemma}[Regularity across the band]\label{lem:band}
In the setting of Lemma~\ref{lem:cutoff},
$\log(\tau_r/\tau_u)\le k^*\log(1+4/\eta)\le\mathsf T_d$.
\end{lemma}
\begin{proof}
Lemma~\ref{lem:transfer}(c) applies, with the $\eta$ of
\eqref{eq:twostep}, because $u+1\le D/2$. By definition
$k^*\le1+\log(2^{77}\Lambda_r/\varepsilon_0)/(2\log(1/\hat\rho))$; use
$\Lambda_r\le d$ (L6) and $\hat\rho\le1.3c_*/\log d$ (L5).
\end{proof}

So $\tau_r/\tau_u\le e^{\mathsf T_d}=d^{o(1)}$, uniformly in $n$, with
$u-r=2k^*=O(\log d/\log(1/\rho))$ for $\rho=r/d$. The convergence is slow:
at $\sigma_1=1$ and $d=10^{176}$, $\mathsf T_d=136.5$, while the leading
term $\frac12\log(1+\frac4\eta)\log d/\log\log d$ is $234.7$. The
difference comes mainly from the term $\log(1/(1.3c_*))\approx6.6$ in the
denominator $2\log(\log d/(1.3c_*))$ of $\mathsf T_d$, which at this $d$
still exceeds $\log\log d\approx6.0$. (These values are computed by the
script of Remark~\ref{rem:d0}.)

\begin{proof}[Proof of Theorem~\ref{thm:S1}]
The statement about $\liminf_NF_N(S)$ follows from
Proposition~\ref{prop:width}(iii), since
$\{q_S\le13/20\}\subseteq\{q_S<2/3\}$.

\emph{Regime I: $n\le n_0(d)$.} Apply Theorem~\ref{thm:signonly} with
$r=t-1$; here $t\ge1$ because $d\ge25$. Lemma~\ref{lem:abfr} gives
$\err_{t-1}(\chi_{[d]})\ge\frac12-\frac t{\sqrt{2\pi d}}
\ge\frac12-\frac1{5\sqrt{2\pi}}>0.4202$, and $C_0n/\binom dt\le1/23<0.0435$.
So $F_\theta(S)>0.3767>3/8$ for every training sequence.

\emph{Regime II: $n_0(d)<n\le n_{\max}(d)$.} Take $r$, $k^*$ and $u$ from
Lemma~\ref{lem:cutoff}, and let $\mathsf E_{\rm coer}=\mathrm{Coer}(r)$ and
$\mathsf E_{\rm gram}=\{\|Q_{r,SS}\|_{\rm op}\le8\,\E\|Q_{r,SS}\|_{\rm op}\}$.
Since $64r\le d$ and $\mathsf B(d,r,n)\le n$,
Theorem~\ref{thm:leancoer} gives $\Pp(\mathsf E_{\rm coer})\ge15/16$, and
Markov's inequality gives $\Pp(\mathsf E_{\rm gram})\ge7/8$. Fix
$S\in\mathsf E_{\rm coer}\cap\mathsf E_{\rm gram}$, an event of probability
at least $1-\frac1{16}-\frac18=\frac{13}{16}$.

\emph{Step 1 (head and band).} By Lemma~\ref{lem:tailgram},
Lemma~\ref{lem:intdim-tanh} (which gives
$\log(\tau_r/\lambda_{r+1})\le\Lambda'_r$) and Lemma~\ref{lem:cutoff}(c),
\[
 4\|Q_{r,SS}\|_{\rm op}+n\lambda_{r+1}\le32\Bigl[(e-1)n\lambda_{r+1}
 +\tau_r\log\frac{\tau_r}{\lambda_{r+1}}\Bigr]+n\lambda_{r+1}
 \le\tau_r\mathsf B_r\le\mathsf B_*d\,\tau_r,
\]
the last step by (L6).

\emph{Step 2 (amplitude).} By Lemma~\ref{lem:band},
\[
 \frac8{c_\theta^2\tau_u}\bigl(4\|Q_{r,SS}\|_{\rm op}+n\lambda_{r+1}\bigr)
 \le\bar{\mathcal R}:=\frac{8\mathsf B_r}{c_\theta^2}\Bigl(1+\frac4\eta
 \Bigr)^{k^*}.
\]
Let $J=\max\{1,\lceil\log(32\bar{\mathcal R})\rceil\}$. Then
$\bar{\mathcal R}e^{-J}\le1/32$. Since $\mathsf B_r\ge\Lambda_r\ge120$ and
$c_\theta<1$, we have $32\bar{\mathcal R}\ge256\mathsf B_r/c_\theta^2>1$, so
$J=\lceil\log(32\bar{\mathcal R})\rceil$, and by Step~1 and
Lemma~\ref{lem:band}, $J\le1+\log(32\bar{\mathcal R})\le\mathsf J_d$.

\emph{Step 3 (admissibility of $J$).} By (L8) and
Lemma~\ref{lem:cutoff}(b), $J\le\mathsf J_d\le e^2(t-4)\le e^2r\le e^2u$.

\emph{Step 4.} Apply Proposition~\ref{prop:master} with these $r$, $u$ and
$J$. By (L8) and Lemma~\ref{lem:cutoff}(d),
$e\sqrt{Ju/(\pi d)}\le e\sqrt{\mathsf J_d\mathsf r_d/\pi}\le1/16$ and
$\frac12e^{-Jd/(8e^2u)}\le\frac12e^{-1/(8e^2\mathsf r_d)}\le1/64$. The
amplitude term is at most $\bar{\mathcal R}e^{-J}\le1/32$, and the last
term is at most $(2+8/c_\theta^2)\varepsilon_0=1/64$. Hence
$F_\theta(S)\ge\frac12-\frac1{16}-\frac1{64}-\frac1{32}-\frac1{64}=\frac38$.
\end{proof}

\begin{remark}[Status, and the role of the coercivity theorem]\label{rem:S1status}
Theorem~\ref{thm:S1} is proved. Regime~I uses only Lemma~\ref{lem:abfr} and
Theorem~\ref{thm:signonly}. Regime~II also uses Theorem~\ref{thm:leancoer},
which this paper does not prove: it is a theorem of the collection's Lean
project, closed with the standard axioms (Section~\ref{sec:mean-train}).
The rest of the proof is conventional mathematics and is not formalized;
Theorem~\ref{thm:S1} itself is not a Lean theorem.

The same proof also runs with \cite[Lemma~4.4]{subcritical} at
$\delta=1/8$ in place of Theorem~\ref{thm:leancoer}: take $r$ largest odd
with $A_{1/8}D_r\log(2D_r)\le n$, put $A_{1/8}$ in place of $2^{77}$ in
$\mathsf B_*$, $\mathsf B_r$, $\mathsf T_d$, (L7) and
Lemma~\ref{lem:cutoff}, put $2\log4+2(r+2)\log\frac{ed}{r+2}$ in place of
$\Lambda_r$, and replace (L1) by
$A_{1/8}D_{t-3}\log(2D_{t-3})\le n_0(d)$. In Step~2 the inequality
$32\bar{\mathcal R}>1$ then follows from $\mathsf B_r\ge A_{1/8}\Lambda_r>1$
and $c_\theta<1$. This route gives success probability
$3/4$, with $d_0$ depending also on $A_{1/8}$, which is not explicit, and
it is conditional on \cite[Lemma~4.1]{subcritical}. Indeed the proof of
\cite[Lemma~4.4]{subcritical} runs through \cite[Lemma~4.3]{subcritical}
to \cite[Lemma~4.1]{subcritical}. That lemma does not apply a stated
theorem of \cite{lee}: it uses the intermediate metric estimate (2.20) in
the proof of \cite[Lemma~2.12]{lee}, with the energy normalization of that
proof undone, and \cite[Remark~4.2]{subcritical} justifies this homogeneous
form. We have not re-checked it against \cite{lee}.

The paper proof of the subcritical theorem \cite[Theorem~1.1]{subcritical}
uses \cite[Lemma~4.4]{subcritical}, but its Lean proof uses
Theorem~\ref{thm:leancoer} instead. So the subcritical theorem, in its Lean
form \texttt{MainProofs.\allowbreak subcriticalExponential}, does not
depend on \cite[Lemma~4.1]{subcritical}. That this Lean statement
expresses \cite[Theorem~1.1]{subcritical} (probability at least $3/4$ and
confidence-failure fraction at least $1/16$, for every $\gamma<1$) is the
finding of the project's semantic review (\texttt{lean/\allowbreak
review/\allowbreak subcritical-\allowbreak main-\allowbreak
semantic-\allowbreak review.txt}), not something Lean checks. The Lean
proof covers only this conclusion. As the project's coverage inventory
(\texttt{lean/\allowbreak review/\allowbreak formal-\allowbreak
coverage-\allowbreak inventory.txt}) notes, it does not cover the stronger
regime theorems of \cite{subcritical}: Gibbs error at least $3/8$ at
$\gamma=0$ \cite[Theorem~9.2]{subcritical}, and $F_{13/20}\ge3/8$ for the
sign posterior at $\gamma<0$ \cite[Theorem~10.2]{subcritical}. Their paper
proofs use \cite[Lemma~4.4]{subcritical}, so they rest on
\cite[Lemma~4.1]{subcritical}. Section~\ref{sec:obstruction}
uses \cite[Lemma~4.4]{subcritical} for general $\delta$ and
Theorem~\ref{thm:leancoer} at $\delta=1/16$
(Proposition~\ref{prop:uncond}).

The log-linear size of the budget is not cosmetic. A contraction argument
gives squared coercivity only for $n\gtrsim D_r^2$. With that sample size the band term
$n\lambda_{r+1}/\tau_r$ would be of order $D_r^2/\binom D{r+2}$, the
amplitude ratio would have logarithm of order $r\log(d/r)$,
Corollary~\ref{cor:restriction} would force $r^2\log(d/r)\lesssim d$, and the
argument would reach only $\log n=O(\sqrt{d\log d})$, inside the range of
Theorem~\ref{thm:signonly}. (This counts the method's own bounds; it is not
a statement about the learner.)
\end{remark}

\begin{remark}[Constants, and the size of $d_0$]\label{rem:d0}
(1) The constant $c_*=10^{-3}$ is not optimized. The budget constants
$1/32$, $1/64$ and $\varepsilon_0$ can be made arbitrarily small at the
cost of additive constants in $J$, and the factor $1.1$ in $\mathsf r_d$
can be brought to $1+o(1)$, so the same proof gives every fixed
$c_*<\pi/(64e^2)=6.6\ldots\cdot10^{-3}$. The method stops there because
$\log\bar{\mathcal R}=(1+o(1))\log d$, so Corollary~\ref{cor:restriction}
forces $\rho_u\lesssim1/\log d$.
(2) Every constant in (L1)--(L8) is explicit, so $d_0$ can be computed,
but the conditions make it astronomically large. For $L=2$, the condition
$\mathsf J_d\mathsf r_d\le\pi/(256e^2)$ of (L8) holds at every scanned $d$
with $\log_{10}d$ above $105.5$, $122.7$, $136.2$, $176.0$, $259.1$,
$387.5$ and $593.8$ for $\sigma_1=8,3,2,1,0.5,0.3,0.2$, and for these gains
all other conditions hold at every scanned $d$ above $10^{28.6}$, the last
of them being (L7). The cause is the slow convergence of
$\mathsf J_d/\log d$ to $1$, which the explicit budget slows further:
$\log\mathsf B_*$ and $\mathsf T_d$ each contain $\log2^{77}\approx53.4$.
Just below the threshold, at $\sigma_1=1$ and $d=10^{176}$,
$\mathsf J_d/\log d=1.5099$ still exceeds the allowed value $1.5098$. The
$A_{1/8}$ route of Remark~\ref{rem:S1status} would give a smaller
threshold for this condition if $A_{1/8}$ were small (for instance
$\log_{10}d\approx58.2$ at $\sigma_1=1$ if $A_{1/8}=1$), but $A_{1/8}$ is
not known.
(3) Even apart from $d_0$, the theorem adds to Theorem~\ref{thm:signonly}
only where $n_{\max}(d)>n_0(d)$, that is, for $d$ above a crossover
$d_\times\approx2.58\cdot10^8$; the ratio $\log n_{\max}(d)/\log n_0(d)$ is
$0.67$, $1.78$ and $4.82$ at $d=10^8,10^9,10^{10}$. A direct evaluation of
the inequalities that Regime~II actually uses, with the actual $r$, $k^*$,
$u$ and $J$ in place of the uniform bounds of (L1)--(L8), found all of them
satisfied at $29$ grid dimensions between $10^{8.42}$ and $10^{300}$, for
$L=2$ and $\sigma_1\in\{0.3,1,3\}$; the largest value of $Ju/d$ was
$1.54\cdot10^{-3}$, at $d=10^{8.42}$ and $\sigma_1=0.3$, against the
requirement $1.66\cdot10^{-3}$. This is a numerical observation at grid
points; it is not part of the theorem. The script
\texttt{numerics/\allowbreak s1\_constants.py} computes the numbers in (2)
and (3) and those in the paragraph after Lemma~\ref{lem:band}; its output
is \texttt{numerics/\allowbreak s1\_constants.txt}.
\end{remark}

\begin{remark}[Why the restriction lemma succeeds where clipped parity does
not]\label{rem:loglinear}
For $h$ of degree at most $u$, the clipped-parity lemma
\cite[Lemma~6.1]{subcritical} gives
$\Pp_x\{fh\ge b\}\le\frac12(1+\E f\clip(h/b))\le\frac12+\frac{\rho_u}2
\mathcal R^{1/2}$, with $\mathcal R$ as in Corollary~\ref{cor:restriction};
the second step uses the Laplacian form \cite[eq.~(23)]{subcritical}, which
does not see the constant term. This bound is linear in the amplitude $\mathcal R^{1/2}$, and
Corollary~\ref{cor:restriction} is logarithmic in it. For the arcsine kernel
the sharp ratio $\theta_r=\binom D{r+2}\lambda_{r+1}/\tau_r\asymp1/r$
\cite[Lemma~3.2]{subcritical} cancels the factor $r$ of the coercivity
threshold, the amplitude ratio is bounded at fixed $\rho=r/d$, and the
linear bound suffices at linear degree, in the paper proof of
\cite[Theorem~10.2]{subcritical}. That proof also uses
\cite[Lemma~4.4]{subcritical}, and the Lean project formalizes only the
weaker conclusion of \cite[Theorem~1.1]{subcritical}
(Remark~\ref{rem:S1status}).
For $\tanh$ the bounds proved here give only an amplitude ratio of order
$\mathsf B_r\,d^{o(1)}$, where $\mathsf B_r$ is of order $rH_r$, with
$H_r=\log(eD/r)$, through both $n/\binom D{r+2}$ (the logarithm in the
coercivity threshold) and $\log(\tau_r/\lambda_{r+1})\ge\log\binom D{r+2}$
(the intrinsic dimension). With the linear bound one would need
$\rho_u(rH_r)^{1/2}d^{o(1)}$ to be small, which confines $r$ to at most
$d^{2/3+o(1)}$. With the logarithmic bound one needs only $\rho_u\log d$ small,
which allows $r$ up to $c\,d/\log d$; this is the source of the exponent
$d\log\log d/\log d$.
\end{remark}

\begin{remark}[Other activations]\label{rem:S1act}
The proof uses $\tanh$ only through three properties: oddness and
Assumption~\ref{ass:act} (for the pairing and the width limit); two-step
regularity of the first-layer coefficients, $c_{m+2}(q_1)\ge\eta c_m(q_1)$
for odd $m$ with some $\eta>0$; and $c_1(q_\ell)>0$ for every layer. These
give \eqref{eq:twostep}, hence $c_m>0$ for every odd $m$ and $K\succ0$. No
lower bound on individual eigenvalues and no far-pair estimate is used. So
Theorem~\ref{thm:S1} holds verbatim, with the same status, for every
activation with these properties, with $d_0$ depending on $\eta$ and $c_1$.
An example is $\operatorname{erf}$. Its normalized dual is
$\widehat\phi_q(\rho)=\arcsin(\varsigma_q\rho)/\arcsin\varsigma_q$ with
$\varsigma_q=2q/(1+2q)$, so
$c_{2k+3}(q)/c_{2k+1}(q)=\varsigma_q^2(2k+1)^2/((2k+2)(2k+3))
\ge\varsigma_q^2/6$ and $c_1(q)=\varsigma_q/\arcsin\varsigma_q>0$. The margin learner for
$\operatorname{erf}$ remains uncovered (Remark~\ref{rem:activations}).
\end{remark}

\begin{remark}[What would give $\exp(\Omega(d))$ for the sign-only
learner]\label{rem:CG}
The following is an \emph{extension suggested by the proof}, conditional on
two open estimates; it is not proved. In the proof of
Theorem~\ref{thm:S1} the amplitude ratio is polynomial in $d$ for two
reasons, the logarithm in the coercivity threshold and the intrinsic
dimension (Remark~\ref{rem:loglinear}). Fix $c_\flat\in(0,c_*)$ and
consider:
\begin{enumerate}
\item[(C)] \emph{Log-free one-sided coercivity at linear degree.} There are
$\rho_0\in(0,1/64]$ and, for each $\delta\in(0,1)$, functions
$\varkappa_\delta,c^{\rm coer}_\delta:(0,\rho_0]\to(0,\infty)$ with
$\rho\log(\varkappa_\delta(\rho)/c^{\rm coer}_\delta(\rho))\to0$ as
$\rho\to0$, such that for all $d\ge d_{\rm C}(\delta,\rho)$, all odd $r$
with $c_\flat d/\log d\le r\le\rho d$, and all $n$ with
$\varkappa_\delta(\rho)D_r\le n\le\varkappa_\delta(\rho)D_{r+2}$,
independent uniform inputs satisfy, with probability at least $1-\delta$,
$\frac1n\sum_i(f(X_i)P(X_i))_-^2\ge c^{\rm coer}_\delta(\rho)\|P\|_2^2$ for
every $P$ of degree at most $r$.
\item[(G)] \emph{Log-free tail Gram bound.} In the same range, with
probability at least $1-\delta$,
$\|Q_{r,SS}\|_{\rm op}\le C^{\rm Gram}_\delta(\tau_r+n\lambda_{r+1})$.
\end{enumerate}
Theorem~\ref{thm:leancoer} has the form of (C) at $\delta=1/16$, with
$c^{\rm coer}=1/4$ and with $2^{76}(\log(8n)+\lceil\log_2D_r\rceil+1)$ in
place of $\varkappa_\delta$, and \cite[Lemma~4.4]{subcritical} has it with
$\varkappa_\delta=A_\delta\log(2D_r)$; both factors grow with $d$, which
(C) does not allow.
Equation~\eqref{eq:gram-lower} shows that (G) would be sharp up to the
constant. Run the proof of Theorem~\ref{thm:S1} with the cutoff
$\varkappa_\delta(\rho_*)D_r\le n<\varkappa_\delta(\rho_*)D_{r+2}$ for a
small fixed $\rho_*$, and with coercivity constant
$c^{\rm coer}_\delta(\rho_*)$ in place of $1/4$ in
Lemma~\ref{lem:meantrain}(c). Since $c_\flat<c_*$, for large $d$ every
cutoff $r$ that this rule assigns to an $n>n_{\max}(d)$ satisfies
$r\ge c_\flat d/\log d$, the lower limit in (C), so the sample sizes
covered by (C) join those of Theorem~\ref{thm:S1} (a lower limit
$d/\log d$ in (C) would leave a gap between the two ranges). Then
$n/\binom D{r+2}\le2\varkappa_\delta(\rho_*)$; (G) bounds the head; a band
of bounded width
$u-r=O(\log(\varkappa_\delta(\rho_*)/\varepsilon_0)/\log(1/\rho_*))$ makes
the last term of Proposition~\ref{prop:master} small; and
Lemma~\ref{lem:transfer}(c) bounds $\tau_r/\tau_u$ independently of $d$.
The amplitude ratio is then bounded independently of $d$, and
Corollary~\ref{cor:restriction} closes for $r\le\rho_*d$ once $\rho_*$ is
small. This would give $\Pp_S\{F_\theta(S)\ge3/8\}\ge3/4$ for all
$n\le e^{cd}$, an $\exp(\Omega(d))$ bound. In the exploration record of
this revision (folder \texttt{verification/\allowbreak
revision-2026-10-20/}: \texttt{exploration-synthesis.md}, Section~2.2, with
the estimate stated in \texttt{lemma-T-draft.md}), (G) is replaced by a
sharp level-mass estimate, $\theta_r\le C/\sqrt r$ together with
$W_i\le CW_{r+2}$ for odd $i\in[r+2,D/2]$, combined with the Tropp-type
bound behind Lemma~\ref{lem:gram}; that variant is sketched there. We do
not prove the level-mass estimate (\texttt{lemma-T-draft.md} is an
unrefereed draft argument), and the computations of
Section~\ref{sec:numerics} do not settle its upper bound
$\theta_r\le C/\sqrt r$: they show $\theta_r\ge c/\sqrt r$ in the computed
range, but $\theta_r\sqrt r$ still grows slowly at its upper end.
Hypothesis~(C) alone does not suffice with the
Gram bounds proved here: Lemma~\ref{lem:tailgram} retains the term
$\tau_r\log(\tau_r/\lambda_{r+1})\ge\tau_r\log\binom D{r+2}$, of order
$d\,\tau_r$ at linear degree, and the bound behind Lemma~\ref{lem:gram}
retains a term polynomial in $d$ unless the level masses are regular in the
above sense.
\end{remark}
